\chapter{Las redes que no están}

\label{cap:redes}
El agua, la cloaca y el gas son el punto donde la conversión de suelo deja de ser una
cuestión de planos. Este capítulo muestra qué proyectos existieron, qué se le exigió a cada
emprendimiento y qué queda hoy: un corredor donde el Estado proyectó redes públicas y
homologó, en los mismos años, abastecimientos privados. Es la forma más concreta en que la
asimetría que enuncia el capítulo \ref{cap:planteo} se vuelve cotidiana.

\section{Lo que hay y lo que falta}

Lo que sigue es el inventario de lo que hay y de lo que falta: qué emprendimientos resolvieron su abastecimiento por cuenta propia, qué obras públicas se proyectaron y en qué estado quedaron.

\subsection*{Un tercer club de campo, y su agua}

\begin{ficha}{Concesión de agua pública --- B.O.\ Nº 19.452 y 19.453, 29 y 30 de diciembre de 2014}
Expediente 34-6147/05. El \textbf{Dr.\ Gonzalo Saravia Echevehere}, administrador judicial en el juicio sucesorio de don \textbf{Adolfo Saravia Valdez}, promotor del loteo \textbf{``Lesser''} a desarrollarse en el \textbf{catastro 1869 del departamento La Caldera}, tiene iniciado trámite de concesión de abastecimiento poblacional a favor de la población que conforma el loteo \textbf{``Finca Lesser''}, ``\textbf{desarrollo urbano realizado bajo la figura de Club de Campo}''.

\textbf{Aguas a captar desde el río Lesser}, en las coordenadas 24° 41' 26,71'' S y 65° 28' 30,00'' O:
\begin{itemize}[leftmargin=*,nosep]
\item \textbf{115 m³ por día} para uso poblacional, \textbf{de ejercicio permanente};
\item un caudal adicional para \textbf{uso recreativo}, \textbf{de ejercicio eventual}.
\end{itemize}

Firma: Dr.\ Rafael Ángel Figueroa, Jefe del Programa Jurídico.
\end{ficha}

Este aviso completa el cuadro que este capítulo venía armando pieza por pieza, y lo completa con una precisión que ninguna de las piezas anteriores permitía.

\textbf{Es el tercer club de campo del departamento} que aparece en el relevamiento, después de Las Vertientes y de Chacras del Gallinato. Con la Urbanización Cerro de Buena Vista, son cuatro desarrollos de baja densidad con trámite documentado entre 2014 y 2015.

\textbf{Y su agua es permanente.} El abastecimiento poblacional de Finca Lesser se capta del río Lesser con carácter de \textbf{ejercicio permanente}, invocando el artículo 66 del Código de Aguas. Compárese con la concesión del Catastro 1700, seis meses anterior: ciento veinte habitantes abastecidos por un pozo perforado, con carácter \textbf{eventual}, porque el artículo 143 declara eventuales todas las concesiones subterráneas sin excepción.

De modo que el mismo Estado, el mismo año y en el mismo departamento tramitó dos derechos de agua para consumo humano con \textbf{jerarquías jurídicas distintas}. La diferencia no está en el destino --- ambos son abastecimiento poblacional --- ni en el número de beneficiarios. Está en la fuente: quien puede captar de un río tiene derecho firme; quien depende de un pozo, no.

Y captar de un río exige lo que un pozo no exige: una fracción con frente sobre el curso, una obra de toma, un estudio y un trámite que en este caso lleva expediente de \textbf{2005}. Es decir, nueve años. \textbf{La permanencia del derecho de agua es, en los hechos, una función de la capacidad de esperar nueve años y de poseer tierra sobre la ribera.}

\begin{cautela}
Dos precauciones sobre este documento.

\textbf{La primera, de lectura.} El caudal consignado para uso recreativo aparece en el aviso con una cifra que, tomada literalmente, resultaría desproporcionada respecto del uso poblacional autorizado en la misma resolución. Es muy probable que se trate de un error de reconocimiento del escaneo o de una unidad mal transcripta. \textbf{Este libro no publica esa cifra} hasta contrastarla contra el original impreso.

\textbf{La segunda, de fondo.} El promotor actúa como \textbf{administrador judicial de una sucesión}. El loteo Lesser se desarrolla, entonces, sobre un acervo hereditario en trámite, y quien gestiona la concesión lo hace por mandato de un juez. Es una figura que este libro no había encontrado hasta ahora y que introduce un actor nuevo en el circuito de conversión de suelo: el juzgado sucesorio. Qué expediente, qué juzgado y con qué autorización se dispuso el desarrollo urbanístico de un bien sucesorio es materia de ese expediente, no de este aviso.
\end{cautela}

\begin{ficha}{Concesión de agua pública --- B.O.\ Nº 19.963 a 19.967, 14 a 21 de febrero de 2017}
Expediente 0090034-263317/2015. El titular registral del \textbf{catastro 4095} del departamento La Caldera gestiona la concesión de uso de agua pública «para abastecimiento de \textbf{630 familias} del Club de Campo \textbf{Las Vertientes--Eco Pueblo}», sito en ese catastro: \textbf{600,48 m³ por día, con carácter permanente}, a derivar de la \textbf{acequia El Alto, margen derecha del río Wierna}. Encuadre: arts.\ 24 inc.\ a), 51, 62, 65, 66, 118 a 125 y 201 del Código de Aguas y Decreto 2299/03. Oposiciones en treinta días hábiles. Firma Silvia F.\ Santamaría como jefa del Programa Fiscalización y Control; edicto del 7 de febrero de 2017.
\end{ficha}

\begin{cautela}
Con esto Las Vertientes, el primero de los clubes de campo de este capítulo, tiene su agua documentada, y la comparación se amplía. Finca Lesser gestionaba 115 metros cúbicos diarios; Las Vertientes, en 2017, más de cinco veces esa cantidad, para 630 familias: unos 950 litros por familia por día. El derecho pedido es permanente y se toma de una acequia del Wierna, el mismo río cuyas márgenes la Provincia determinaba en esos meses (capítulo \ref{cap:ribera}). Y la solicitud es de 2017, tres años después de que Inmuebles aprobara el club con la condición de un «reglamento de condominio con servicio de agua privado» (capítulo \ref{cap:loteo}): el agua que el plano presuponía privada se gestiona después como concesión pública. Cuatro años más tarde, en 2021, los vecinos de la zona alta de Vaqueros reclamaban agua de camión. El aviso no dice si la concesión se otorgó ni con qué caudal. \pendiente{resolución que otorgó o denegó la concesión del expediente 0090034-263317/2015, y relación entre el catastro 4095 y la matrícula 3989 del plano de 2014}
\end{cautela}

\subsection*{Quién firma, once años después}

Las nóminas de autoridades publicadas a fines de 2025 ---la de la estructura de la Jefatura de Gabinete, donde la Decisión Administrativa 569/25 ubica a la Subsecretaría de Recursos Hídricos (capítulo \ref{cap:ribera})--- permiten cerrar un círculo y abrir dos puertas.

\begin{ficha}{Subsecretaría de Recursos Hídricos --- estructura vigente al 30/12/2025}
\begin{itemize}[leftmargin=*,nosep]
\item Coordinador: David Le Favi.
\item \textbf{Dirección General de Cuencas y Política Hídrica}: Jorge Genaro Torres.
\item \textbf{Dirección General de Concesiones Hídricas}: Leandro Bruno.
\item \textbf{Dirección Legal y Técnica: Silvia Fabiana Santamaría\index{Santamaría, Silvia F.}}.
\end{itemize}
En la misma nómina, la \textbf{Secretaría de Tierra y Bienes del Estado} registra a Miriam Esther Ochoa como Directora de Regularización de Tierras Fiscales, Rurales y Urbanas, y a Natalia Macarena Pierola Matos como Directora General de Regularización Dominial.
\end{ficha}

\begin{cautela}
La continuidad es el dato. \textbf{Silvia F.\ Santamaría\index{Santamaría, Silvia F.}} es quien firmó, como Abogada Jefe del Programa Fiscalización y Control, el edicto de junio de 2014 por el que se tramitó la concesión de agua de pozo del Catastro 1700 de La Caldera. Once años y medio después es Directora Legal y Técnica del mismo organismo. Entre una fecha y otra ocurrieron los dos aluviones, el amparo, la condena, la contratación de la consultoría y la aprobación del Plan.

Eso no es un señalamiento personal y este libro no lo formula como tal. Es lo contrario de lo que suele suponerse sobre la administración provincial: en el área hídrica hay memoria institucional disponible, radicada en personas identificables que siguen en funciones. Quien quiera preguntar por qué la línea de ribera se fijó en 2014 y el desborde ocurrió en 2018 no tiene que reconstruir un organismo desaparecido: tiene que escribirle a una dirección que existe.

Y las dos puertas que la nómina abre son literales. La \textbf{Dirección General de Concesiones Hídricas} es la oficina exacta donde vive la nómina de concesiones del departamento que este capítulo pide. La \textbf{Dirección General de Cuencas y Política Hídrica} es donde vive el Plan de Manejo y la planificación de obras que la jueza ordenó presentar por plazos. Ninguno de los dos pedidos tenía hasta ahora destinatario nombrado.
\end{cautela}

\begin{cautela}
Un cambio de rango que corresponde consignar con reservas. Los documentos que este libro cita de 2014, 2021 y 2024 nombran al organismo como \textbf{Secretaría} de Recursos Hídricos, y en 2024 lo ubican bajo el Ministerio de Producción y Desarrollo Sustentable. La nómina de diciembre de 2025 lo nombra como \textbf{Subsecretaría} de Recursos Hídricos, dentro de la Jefatura de Gabinete, aunque la Ley 8.511 asigna los recursos hídricos al Ministerio de Producción y Minería (capítulo \ref{cap:ribera}).

Los edictos del propio organismo fechan el cambio: en julio de 2025 firma la Secretaría y en marzo de 2026 la Subsecretaría (capítulo \ref{cap:ribera}). El descenso de rango está, por lo tanto, acreditado. Lo que el organigrama oficial y la ley de ministerios deben aclarar es si arrastró competencias, presupuesto o planta, y eso sigue siendo pregunta.
\end{cautela}

\subsection*{La región metropolitana existe, pero solo para los colectivos}

Hay una figura legal que reconoce formalmente lo que este libro describe --- que el corredor norte funciona como periferia de la capital --- y que conviene poner sobre la mesa por lo que incluye y por lo que deja afuera.

\begin{ficha}{Región Metropolitana de Salta --- Ley 7.322, en el marco de la Ley 6.994}
Establece la Región Metropolitana de Salta respecto de los \textbf{servicios de transporte por automotor de pasajeros}, propios e impropios, urbanos e interurbanos. La integran los municipios de \textbf{Salta, San Lorenzo, Vaqueros, Cerrillos, Rosario de Lerma, Campo Quijano, La Merced y La Caldera}, más los que el Poder Ejecutivo incorpore en el futuro.

Esos servicios pasan a ser \textbf{competencia provincial}. La \textbf{Autoridad Metropolitana de Transporte} es el organismo de aplicación, con potestades de planificación, organización, regulación, fiscalización y control, y con facultad de dictar los reglamentos sobre recorridos, cupos, permisos, regímenes tarifarios y calidad del servicio.

\textit{Fuente: considerandos de la Resolución A.M.T.\ Nº 376/14, B.O.\ Nº 19.315.}
\end{ficha}

Léase la lista de municipios: es exactamente el conjunto que este libro viene tratando como un área funcional única. El Estado provincial la reconoció, le puso nombre, la sacó de la órbita municipal y le creó una autoridad con poder de policía.

\textbf{Lo hizo para el transporte de pasajeros.} \textit{(Con una salvedad que el apartado sobre el CorInDeS, más abajo, desarrolla: en 2009 los cinco municipios del corredor se constituyeron por convenio como microrregión, con el ordenamiento territorial entre sus fines. Lo que sigue vale para la única figura con ley, organismo y presupuesto.)}

No para el agua, que en Vaqueros se proyecta por localidad y en los clubes de campo se resuelve en forma privada. No para el suelo, cuya aprobación de loteos tramita municipio por municipio. No para el río Vaqueros, que separa dos jurisdicciones y que ninguna de las dos administra en su conjunto. No para los áridos, que se conceden a ambos lados del límite departamental sin autoridad común.

Cuando este libro sostiene que el corredor se gobierna en fragmentos, no está describiendo una imposibilidad institucional ni una laguna del federalismo. \textbf{La provincia sabe construir autoridades metropolitanas cuando decide hacerlo}, y lo demostró con los colectivos. Que no exista una figura equivalente para la cuenca, para el suelo o para el agua no es un problema de diseño constitucional: es una decisión sobre qué se administra en escala metropolitana y qué se deja en escala municipal.

Un dato que aparece de paso en esos mismos considerandos y que conviene guardar. Al evaluar un reclamo del intendente de San Lorenzo, la Gerencia de Transporte del organismo señala que \textbf{no cuenta con datos del último Censo Nacional por localidad, sólo por departamento}, y estima el crecimiento comparando departamentos: Cerrillos habría crecido un 36 por ciento y \textbf{La Caldera un 40}.

La cifra es del orden de la que el capítulo \ref{cap:pdua} consigna para el período 2001--2010. Lo interesante es lo otro: en 2014, un organismo provincial con potestades de planificación reconocía por escrito que \textbf{planificaba sin datos desagregados por localidad}. Es la misma carencia que este libro tuvo que resolver procesando Redatam por su cuenta.

\subsection*{Vaqueros tuvo su obra de agua en 2014}

\begin{ficha}{Resolución SOP Nº 220 del 21 de abril de 2014 --- B.O.\ Nº 19.304 y 19.306}
Expediente 125-120.154/13. Aprueba el proyecto ejecutivo \textbf{confeccionado por la Compañía Salteña de Agua y Saneamiento} y completado por la Dirección de Obras Hídricas y Saneamiento, para la obra \textbf{``Acueducto, Mejoras en Planta Potabilizadora, Canal de Aducción y Redes Maestras --- Localidad Vaqueros''}.

Presupuesto oficial \textbf{\$1.500.000}, modalidad ajuste alzado, plazo de \textbf{ciento ochenta días corridos}. Se autoriza el proceso selectivo con encuadre en el art.\ 10 de la Ley 6838. En el mismo mes se llama a Concurso de Precios Nº 16/2014.
\end{ficha}

\begin{cautela}
Póngase esta obra al lado del reclamo vecinal que este capítulo recoge.

En abril de 2014 la Provincia aprobó el proyecto ejecutivo de un acueducto con planta potabilizadora, canal de aducción y redes maestras para Vaqueros, con un plazo de ejecución de seis meses. En julio de 2021 --- siete años después --- vecinos de la zona alta del mismo municipio se manifestaban tras más de quince años de reclamos porque en sus casas no había agua.

\textbf{Las dos cosas pueden ser compatibles.} Una obra de redes maestras y planta potabilizadora no necesariamente alcanza la zona alta; un acueducto puede mejorar la conducción sin extender la cobertura; y un proyecto ejecutivo aprobado no equivale a una obra terminada. Este libro no sabe si la obra se ejecutó, ni hasta dónde llegó.

Sí sabe que en noviembre de 2023 la zona alta todavía se estaba conectando por etapas. La \textbf{Resolución Nº 855/23 de la Secretaría de Obras Públicas} (B.O.\ Nº 21.605, 5 de diciembre de 2023; expte.\ 125-263.672/21) aprueba la «Optimización red de agua potable --- Vaqueros --- Zona Alta --- Sector III --- Etapa 3»: 1.686 metros de cañería y \textbf{cincuenta conexiones domiciliarias}, con proyecto de Aguas del Norte, \$39.900.254,77 y sesenta días de plazo, adjudicada por contratación abreviada a la propia Municipalidad de Vaqueros. Sus considerandos dicen que esa tercera etapa necesita, para funcionar, que el IPV concluya «el nuevo acueducto en la zona alta» y que el agua de «la Nueva Planta Potabilizadora de Campo Alegre» llegue a una cisterna de 600 m³ en rehabilitación. \textbf{Nueve años después del proyecto de 2014, la red de la zona alta dependía todavía de obras no terminadas.} Y la serie sigue, con la misma forma que las defensas del río (capítulo \ref{cap:defensas}): en enero de 2025 la Resolución 10/25 aprueba el legajo de la «Optimización de redes Vaqueros --- Zona Alta --- Sector 4 --- Etapa I»; en julio, la 367/25 aprueba un convenio de rescisión de la obra «Optimización de redes Vaqueros»; en septiembre, la 482/25 vuelve a aprobar el legajo del Sector 4, Etapa I. \textit{(Esos tres actos se citan por su título en el índice del Boletín; sus textos no se consultaron.)} En 2026, en cambio, la obra se licita a empresas: el «Nexo para abastecimiento de agua potable a Villa Sara y alrededores y redes distribuidoras», con presupuesto oficial de \$744.570.259,64 y ciento ochenta días de plazo, recibe quince ofertas y se adjudica a NECA S.R.L.\ por \$677.541.957,02 (Resolución 404/26, B.O.\ Nº 22.236, 29 de julio de 2026). Entre las ocho desestimadas está la de D.E.S.A.\ Defensa y Encauzamiento S.A., la misma que en 2022 ofertó por el puente sobre el río Vaqueros, por no presentar el certificado de habilitación de su asesor técnico.

Lo que en cambio puede afirmar, y es suficiente para el argumento del capítulo, es que \textbf{la ausencia de red en Vaqueros no se explica por ausencia de proyectos}. Hubo proyecto ejecutivo, presupuesto asignado y llamado a concurso, en el mismo año en que el Estado exigía a un club de campo del mismo municipio que previera resolver el agua \textbf{en forma privada}.

Ese es el hallazgo: en 2014, sobre el mismo territorio y con meses de diferencia, el Estado proyectaba red pública para la localidad y homologaba abastecimiento privado para el emprendimiento. No son políticas contradictorias --- son dos respuestas a dos capacidades de pago ---, y esa simultaneidad es exactamente lo que este libro llama circuitos paralelos.
\end{cautela}

Hubo, al menos, actividad operativa sobre el sistema. En mayo de 2014 la Compañía Salteña de Agua y Saneamiento contrató directamente, por el artículo 13 h de la Ley 6838, el \textbf{alquiler de una retroexcavadora para trabajos en el Acueducto Norte, localidad de Vaqueros}: proveedor Geomix S.R.L., \$10.100 más IVA, expediente 5.836/14. En el mismo mes adquirió rodamientos para bombeo con destino a la localidad de La Caldera, expediente 5.438/14.

Son montos menores y contrataciones de rutina. Se consignan porque acreditan que en 2014 había actividad operativa sobre el acueducto norte en ambos municipios, y porque ese mismo acueducto reaparece en 2019 en el escrito del amparo, a propósito de un desmonte y de una toma de agua no autorizada.

\section{Agua, cloaca y gas: lo que se proyectó y lo que falta}

\subsection*{La única obra de red de agua registrada}

\begin{ficha}{Una obra que sí existe: red de agua de La Calderilla}Registro del sistema nacional de seguimiento de inversión pública, consultado en septiembre de 2026:
\begin{itemize}[leftmargin=*,nosep]
\item \textbf{Denominación:} ``Optimización del servicio de La Calderilla, La Caldera, Salta''. Renovación completa de las redes y de las conexiones domiciliarias de La Calderilla.
\item \textbf{Organismo responsable:} Ente Nacional de Obras Hídricas de Saneamiento (ENOHSA). Jurisdicción beneficiaria: Secretaría de Obras Públicas de Salta. Modalidad: convenio. Código BAPIN 132888.
\item \textbf{Costo estimado:} \$114.372.451.
\item \textbf{Inicio:} 2021. \textbf{Finalización prevista:} 2023.
\item \textbf{Estado:} \textbf{en ejecución. Avance: 47,43 por ciento.}
\end{itemize}
\end{ficha}

Corresponde leer el cuadro por lo que dice y no por lo que sugiere. Dice que la única obra de red de agua registrada para el departamento en el sistema nacional tenía fecha de finalización en 2023 y, tres años después de esa fecha, figura en ejecución con menos de la mitad de avance.

No dice por qué, pero el Boletín provincial da al menos una causa: en octubre de 2024 el Ministerio de Infraestructura aprobó la \textbf{rescisión del contrato} con Norte Áridos S.R.L.\ para la obra «Optimización del servicio de La Calderilla --- nueva red distribuidora» (Resolución 154/24, B.O.\ Nº 21.828; capítulo \ref{cap:cot}). Si el 47,43 por ciento es el avance que quedó congelado con esa rescisión, o si la obra siguió con otro contratista, el registro no lo distingue; tampoco puede descartarse que el porcentaje esté desactualizado, aunque el sitio declara actualización diaria.

Lo que el dato aporta al argumento de este capítulo es de otro orden y no depende de la causa: \textbf{el municipio no controla el ritmo de la única infraestructura de agua que se está construyendo en su territorio}. La obra es de un ente nacional, ejecutada por convenio con la Provincia, sobre una localidad del departamento. La competencia sobre el suelo es municipal; la capacidad de dotarlo de servicio, no.

Ese desacople --- competencia local, capacidad extra-local --- es el mismo que este libro describe para la aprobación de loteos, para el ordenamiento de bosques y para la expansión de la red de gas. Aquí aparece con un número y una fecha de vencimiento.

\subsection*{Las cloacas}

\textbf{Cloacas.} En la audiencia pública por el loteo Lares de la Inmaculada, en abril de 2019, los vecinos plantearon que temían contaminación por la cantidad de familias a instalarse \textbf{en una zona sin cloacas}. El proyecto de ordenanza del Plan de Desarrollo Urbano Ambiental de La Caldera ubica la planta de tratamiento de efluentes, junto con los basurales, entre los equipamientos ``molestos o incompatibles'' que se localizan en zona sub-urbana; y su lista de proyectos complementarios incluye, como pendiente, la conexión de infraestructura hacia la planta depuradora.

\begin{ficha}{Licitación Pública Nº 29/2015 --- B.O.\ Nº 19.532, 5 de mayo de 2015}
Aguas del Norte llama a licitación, en el marco del \textbf{Plan de Obras de Infraestructura Básica de Agua Potable y Saneamiento}, para la obra \textbf{``Redes Colectoras, Colectora Máxima y Planta Depuradora Cloacal --- Vaqueros''}. Expediente 7807/15. Plazo de ejecución: \textbf{ocho meses}. Pliego: \$10.000. Apertura de ofertas: 9 de junio de 2015.
\end{ficha}

\begin{cautela}
La pregunta por el estado de la red cloacal y los proyectos existentes queda parcialmente respondida, y la respuesta es que el proyecto existió y fue completo.

No se trata de una extensión ni de un tramo: la licitación abarca \textbf{redes colectoras, colectora máxima y planta depuradora}. Es un sistema cloacal entero para la localidad, licitado en mayo de 2015, con ocho meses de plazo.

Lo que este libro \textbf{no sabe} es qué ocurrió después: si se adjudicó, si se ejecutó, si se terminó y a qué superficie del ejido alcanzó. Y sabe, en cambio, que cuatro años más tarde los vecinos de un loteo de Vaqueros seguían planteando en audiencia pública el temor a la contaminación ``en una zona sin cloacas''.

Las dos cosas pueden ser compatibles: una planta depuradora con colectora máxima puede servir al casco urbano y no alcanzar a los loteos de la periferia, que son justamente los que crecieron. \textbf{Pero eso es exactamente lo que hay que averiguar}, y con el número de expediente ya no es una búsqueda difusa.
\end{cautela}

\begin{cautela}
Y una circular aclaratoria del 5 de junio de 2015 revela por qué esa obra podía ser más difícil de lo que su pliego sugiere.

Ante consultas de los adquirentes del pliego, el directorio de la empresa aclaró que el ítem ``trámites para servidumbre de paso y acueducto de la obra completa'' debe ser cotizado por los oferentes y ejecutado por la contratista ``en lo que respecta a todos los tramos de la obra completa \textbf{que hayan sido proyectados en terrenos privados}''. Esos trabajos comprenden \textbf{el estudio general de títulos de los terrenos afectados}, la elaboración de los planos de mensura para servidumbre, su presentación y aprobación ante la \textbf{Dirección General de Inmuebles}, y su entrega al comitente.

\textbf{La red cloacal de Vaqueros debía atravesar propiedad privada.} No un tramo puntual: la aclaración habla de ``todos los tramos'' de los cuatro rubros que estén proyectados sobre terrenos privados, lo que sugiere que no eran pocos.

Tres consecuencias.

\textbf{Primera:} el Estado \textbf{transfirió al contratista} la tarea de averiguar quién es dueño de qué. El estudio de títulos --- exactamente el trabajo que este libro no puede hacer porque es arancelado y presencial --- quedó a cargo de la empresa que gane la licitación, y su costo, incorporado al precio de la obra pública.

\textbf{Segunda:} el trazado de la red no lo define la ingeniería sola. Lo define también dónde se consiguen servidumbres. Una colectora máxima que debe negociar paso por fincas privadas termina siguiendo el mapa de la propiedad, no el de la población.

\textbf{Y tercera, que es una hipótesis y como tal se enuncia:} si la obra se ejecutó y no alcanzó a la periferia, la explicación puede no ser presupuestaria ni técnica. Puede estar en las servidumbres. \textbf{Este libro no lo afirma}, porque no sabe si la obra se hizo ni por dónde pasó. Lo señala porque, de confirmarse, sería la demostración más directa de su tesis: la infraestructura pública llega tarde a este corredor porque el suelo se privatizó antes, y ahora hay que pedirle permiso.
\end{cautela}

\subsection*{El gas}

Primero, dónde llega. El Ente Nacional Regulador del Gas publica tres cartografías que permiten responderlo, y las tres dicen lo mismo.

\begin{ficha}{Tres cartografías del ENARGAS sobre el corredor}
\begin{itemize}[leftmargin=*,nosep]
\item \textbf{Atlas provincial}, lámina «Sistema de Transporte y Distribución de Gas Natural de la Provincia de Salta», sin fecha (lámina \ref{lam:gasatlas}). \textbf{Vaqueros figura como ciudad abastecida con gas natural}, en el extremo de un ramal de la malla regional del Valle de Lerma. \textbf{La Caldera no figura}, y al norte de Vaqueros no hay localidad abastecida ni gasoducto hasta el límite con Jujuy. Distribuidora de toda la zona: Naturgy NOA S.A.
\item \textbf{Mapa nacional «Localidades abastecidas con gas distribuido por redes»}, fechado en octubre de 2024 (lámina \ref{lam:gasnacional}). Registra a Vaqueros entre las localidades abastecidas con gas natural a cargo de la distribuidora. La Caldera no aparece.
\item \textbf{Mapa interactivo del ENARGAS}, exportado el 16 de septiembre de 2026 a escala 1:144.448 (lámina \ref{lam:gasvisor}). Muestra un único marcador en el corredor, sobre el área urbana del sur, atravesada por la Ruta 9; ninguno sobre el pueblo de La Caldera ni sobre el embalse.
\item \textbf{Advertencia de los propios mapas:} a esas escalas no se distinguen los barrios dentro de cada localidad abastecida.
\end{itemize}
\end{ficha}

\begin{figure}[p]
\centering
\includegraphics[width=\textwidth,height=0.36\textheight,keepaspectratio]{lam-gas-atlas.jpg}%
\caption[La malla de gas del Valle de Lerma y su extremo norte en Vaqueros]{\textbf{Sistema de transporte y distribución de gas natural, sector del Valle de Lerma.} Los puntos naranja son ciudades abastecidas con gas natural; en violeta, los gasoductos regionales; en azul, el Gasoducto Norte de Transportadora de Gas del Norte S.A.; la línea gris gruesa, al norte, es el límite con Jujuy. \textbf{Vaqueros} es el punto más septentrional de la malla. Entre Vaqueros y ese límite no hay ningún punto: allí está el municipio de La Caldera.}\label{lam:gasatlas}%
\par\vspace{0.3em}%
{\sffamily\scriptsize\color{gris}\raggedright\textbf{Procedencia.}\ ENARGAS, \textit{Atlas provincial del gas}, lámina de la Provincia de Salta, sin fecha, consultada en septiembre de 2026. Recorte del sector central; no se alteró el mapa.\par}%
\vspace{1.4em}
\includegraphics[width=\textwidth,height=0.36\textheight,keepaspectratio]{lam-gas-nacional.jpg}%
\caption[Localidades abastecidas con gas por redes, noroeste, octubre de 2024]{\textbf{Localidades abastecidas con gas distribuido por redes, sector de Salta y Jujuy.} Los círculos rojos llenos son localidades abastecidas con gas natural a cargo de la distribuidora; los partidos, localidades a abastecer con proyectos en ejecución. \textbf{Vaqueros} figura con círculo lleno, junto a San Lorenzo y Salta. \textbf{La Caldera no figura}, ni como abastecida ni como proyecto.}\label{lam:gasnacional}%
\par\vspace{0.3em}%
{\sffamily\scriptsize\color{gris}\raggedright\textbf{Procedencia.}\ ENARGAS, \textit{Localidades abastecidas con gas distribuido por redes. República Argentina}, lámina fechada «ENARGAS -- 10/2024», consultada en septiembre de 2026. Recorte del sector noroeste; no se alteró el mapa.\par}%
\end{figure}

\lamina{lam-gas-visor.jpg}{lam:gasvisor}%
{El corredor en el mapa interactivo del ENARGAS, septiembre de 2026}%
{\textbf{Mapa interactivo del ENARGAS sobre base del IGN, escala 1:144.448.} Arriba, el embalse Campo Alegre; en el centro, sobre el río, el área urbana con el símbolo de cabecera departamental, que corresponde al pueblo de La Caldera; abajo, atravesada por la Ruta 9 y junto al río Mojotoro, el área urbana de Vaqueros. \textbf{El único marcador azul del corredor está sobre Vaqueros.} La exportación no trae leyenda: que el marcador corresponda a la capa de localidades abastecidas es una lectura de este libro, coherente con las dos láminas anteriores, y no una indicación del mapa.}%
{Exportación en PDF del mapa interactivo del ENARGAS (\textit{sig.enargas.gov.ar}), 16 de septiembre de 2026, 7:15. Se recortó el marco del mapa con su fecha y escala; no se alteró el contenido.}

El departamento tiene, entonces, \textbf{un municipio con gas de red y otro sin él}, separados por el río Wierna (Ley 4987; capítulo \ref{cap:cot}). Ninguno de los tres mapas dice hasta dónde llega la red dentro de Vaqueros: si alcanza a los loteos de la periferia y a los clubes de campo, o sólo al casco. Eso se pregunta a la distribuidora y al municipio, y el apéndice \ref{cap:pedidos} lo registra. Lo que sí dicen es que \textbf{la red se detiene en el municipio más cercano a la capital} y no continúa hacia el pueblo que da nombre al departamento.

\textbf{Gas.} El formulario de inscripción para lotes de interés social --- Anexo II de la Resolución 7/16 --- pregunta al solicitante si su vivienda tiene gas y, en caso afirmativo, si es \textbf{de red} o de \textbf{tubo o garrafa}. El Plan Provincial ``Mi Lote'', creado por el Decreto 399/20, promete entregar lotes ``dotados de los servicios públicos esenciales'', y enumera: \textbf{agua, luz eléctrica y apertura de calles}. El gas de red no integra esa enumeración, y las cloacas tampoco.

La enumeración del Decreto 399/20 es el documento más claro sobre el estándar que el Estado provincial se fija a sí mismo. Su programa insignia de urbanización social define los servicios públicos esenciales como agua, luz y calle. Un lote con esas tres cosas cumple la promesa.

Y el formulario de la Resolución 7/16 hace lo mismo que el capítulo \ref{cap:tierrafiscal} documenta para la zona inundable: \textbf{pregunta por el gas de red y el baremo de la Resolución 8/16 no le asigna puntaje}. El Estado registra el dato y no lo usa para decidir.

\subsection*{La ausencia de red no es una imposibilidad legal}

Conviene consignarlo porque cambia el sentido de la carencia.

La Resolución ENARGAS Nº I-910/09, del 9 de octubre de 2009, rigió durante diecisiete años el procedimiento para la expansión de sistemas de distribución de gas. Bajo ese régimen, una obra de ampliación de red podía tramitarse a instancia de la distribuidora, de un subdistribuidor o de \textbf{terceros interesados}, incluidos los futuros usuarios. La metodología definía el Valor de Negocio del proyecto: \textbf{cuando ese valor superaba el monto de la inversión, la obra correspondía íntegramente a la distribuidora}; cuando resultaba inferior, la distribuidora aportaba ese monto y la diferencia la solventaban los usuarios o los terceros interesados, que recibían después una contraprestación --- para los usuarios residenciales, según la Resolución 630/2019, exclusivamente en metros cúbicos de gas.

La distribuidora licenciada para Salta es Naturgy NOA S.A., ex Gasnor S.A., con licencia otorgada por Decreto Nº 2452/92 sobre las provincias de Salta, Jujuy, Tucumán y Santiago del Estero.

\subsection*{El régimen cambió en abril de 2026}

Este apartado se escribió sobre la Resolución I-910/09 y hay que corregirlo, porque esa norma ya no está vigente y porque la corrección abre una vía documental que el libro no tenía.

\begin{ficha}{La secuencia}
\begin{itemize}[leftmargin=*,nosep]
\item \textbf{Res.\ ENARGAS 778/2025}, del 20 de octubre de 2025 (B.O.\ 21/10/2025). No modificó el régimen: \textbf{puso en consulta pública} un proyecto de reglamentación del artículo 16 de la Ley 24.076 --- T.O.\ 2025, por quince días hábiles.
\item \textbf{Res.\ ENARGAS 875/2025}, del 6 de noviembre de 2025. Amplió el plazo por diez días hábiles más, ante requerimientos de los interesados.
\item \textbf{Res.\ ENARGAS 435/2026}, del 24 de abril de 2026 (B.O.\ 27/04/2026). Su artículo 1º aprueba la nueva \textit{Reglamentación de las autorizaciones del Artículo 16 de la Ley Nº 24.076 --- T.O.\ 2025}. Su artículo 2º \textbf{deja sin efecto la Resolución I-910/09 y toda su normativa concordante y complementaria}. Rige desde su publicación; los trámites iniciados antes continúan bajo el régimen anterior.
\end{itemize}
\end{ficha}

La propia resolución enumera tres cambios respecto del régimen anterior: reduce el horizonte de valuación de \textbf{treinta y cinco a diez años}; reemplaza los costos medios derivados de los márgenes tarifarios por costos marginales surgidos de declaraciones juradas, actualizados por una fórmula de mitad índice de precios al consumidor y mitad índice de precios internos al por mayor; y sustituye el flujo de fondos que cada prestadora elaboraba sin estándar por un aplicativo web de uso obligatorio desde la instancia de factibilidad.

\begin{cautela}
El primero de esos cambios importa directamente al corredor. Un horizonte de treinta y cinco años permitía computar como ingreso la demanda de usuarios que todavía no existen; diez años obliga a computar la demanda que madura en una década. Para un tejido de baja densidad, con lotes vendidos y no edificados y con casi una de cada cinco viviendas sin residentes permanentes, ese acortamiento \textbf{endurece el cálculo}: menos consumo computable sobre la misma longitud de cañería.

Dicho de otro modo, el régimen que reemplaza a la I-910/09 hace más difícil, no más fácil, que una expansión hacia La Caldera resulte a cargo de la distribuidora. El libro no puede afirmar que ése haya sido el propósito de la norma, que es de alcance nacional y responde a otra discusión. Puede afirmar que ése es su efecto previsible sobre un parque habitacional con las características que los capítulos anteriores midieron.
\end{cautela}

Y hay un hallazgo documental de otro orden. Los considerandos de la Resolución 435/2026 enumeran a quienes presentaron observaciones durante la consulta pública. Entre ellos figura \textbf{Naturgy NOA S.A.}, con presentación del 27 de noviembre de 2025 en el expediente EX-2025-112087756-APN-GD\#ENARGAS.

Es decir: la distribuidora licenciada para Salta \textbf{tiene posición formal, escrita y pública} sobre la regla que determina quién paga la extensión de red. Este libro venía pidiendo una consulta a la empresa que la empresa no está obligada a contestar. El expediente de la consulta pública es una vía de acceso a esa misma posición que no depende de su voluntad de responder.

De modo que la ausencia de red de gas en el municipio de La Caldera no responde a un impedimento normativo ni a una imposibilidad técnica declarada. Responde a que, hasta donde alcanza este relevamiento, \textbf{nadie solicitó que se calculara el Valor de Negocio del proyecto}. Y ese cálculo es el que determinaría si la obra le corresponde a la distribuidora o si requiere aporte de terceros.

La diferencia entre ``no hay red porque no se puede'' y ``no hay red porque nadie pidió que se evaluara'' es la diferencia entre una carencia y una omisión. Establecer cuál de las dos es requiere consultar al ENARGAS si existen expedientes de expansión para el departamento, y a la distribuidora si evaluó alguna. Ninguna de las dos consultas fue realizada por este libro.

\subsection*{Quién se beneficia de la ausencia}

\begin{cautela}
Corresponde formular la pregunta, y responderla con el mismo criterio que el resto del libro: identificando posiciones estructurales y no intenciones.

\textbf{El desarrollador.} Sin red pública, cada emprendimiento resuelve sus servicios por cuenta propia, y eso deja de ser un costo para convertirse en argumento de venta. Un loteo con vertiente propia vende algo que sus vecinos no tienen. Con red, ese diferencial desaparece. Y hay un efecto adicional sobre el control: el artículo 144 del Código exige garantía de ejecución de obras de urbanización antes de comercializar, pero sin un estándar de red contra el cual medir, ``obras de infraestructura'' es una categoría cuyo cumplimiento nadie puede verificar desde afuera.

\textbf{El propietario no residente de parcela sin edificar.} Es, según el capítulo \ref{cap:poblacion}, un universo identificable, que el municipio intimó en agosto de 2024 por morosidad, y que representa una proporción significativa de un parque con 19,80 por ciento de viviendas sin residentes. Para quien tiene un lote que no habita, la extensión de red no es un beneficio inmediato: es contribución por mejoras y es exposición al gravamen sobre lotes no edificados que ordena el artículo 30 del Código.

\textbf{Nadie decidió esto.} No hay en el material relevado indicio alguno de una decisión de impedir la infraestructura, y el libro no la insinúa. Lo que hay es una configuración en la que la extensión de red no le conviene, en el corto plazo, a ninguno de los actores con capacidad de gestionarla, y le conviene mucho a quienes no la tienen.

Ése es, otra vez, el mismo mecanismo que el capítulo \ref{cap:planteo} enunció: sin red, el único modo de tener servicios es comprar en un emprendimiento que los provea privadamente. El servicio público se convierte en prestación privada y en atributo comercial, y quien no puede pagarlo queda con garrafa, leña y camión municipal dos o tres veces por semana.
\end{cautela}

\section{Agua pública y agua privada}

Este capítulo y el \ref{cap:loteo} discuten qué significa que una concesión sea permanente o eventual. Un
decreto de 1939 introduce una distinción anterior a esa, y más básica: la que separa el agua
sobre la que el Estado decide del agua sobre la que no.

\begin{ficha}{Decreto Nº 2469 del 26 de enero de 1939}
Expediente 632-O. El \textbf{Dr.\ Lucio Ortiz} solicita inscribir en los registros oficiales un
documento firmado por \textbf{Augusto Regis\index{Regis, Augusto}}, Presidente de la Comisión Municipal de La
Caldera, que certifica que la provisión de agua para riego y usos domésticos de la finca
\textbf{«Abra de Lesser»} se hace del \textbf{arroyo del «Abra»}.

\textbf{Se rechaza.} Considerando que, según el propio recurrente, ese arroyo «nace y muere
dentro de la finca» de su propiedad, «no corresponde la inscripción en el Registro Oficial,
ya que \textbf{se trataría de un derecho de uso de agua de dominio privado}» --- y el registro
sólo procede sobre concesiones de agua de \textbf{dominio público}.

Firman Luis Patrón Costas y Carlos Gómez Rincón.
\end{ficha}

\textbf{Primera: hay agua sobre la que el Estado no lleva registro.} Un curso que nace y muere
dentro de un mismo inmueble es, en ese régimen, del dueño del inmueble, y su uso no se
inscribe. De modo que la serie de concesiones que este libro reconstruye --- Urquiza\index{Urquiza, Francisco S.} en 1924,
el pedido denegado de Campo Alegre en 1938, el Catastro 1700 en 2014, El Durazno en 2015, La Misión III en 2015,
Finca Lesser --- \textbf{no es el universo del agua usada en el departamento: es el universo
del agua registrada}. Cuánta se usó sin registro, porque nacía y moría dentro de una finca,
es algo que ningún archivo público puede decir.

Y eso vuelve más pertinente la observación del capítulo \ref{cap:siglo}: si en 1924 el Estado reconoció por
escrito que no había aforado ningún río de la provincia, y si además había cursos enteros
fuera de todo registro, entonces \textbf{nadie estaba en condiciones de saber cuánta agua
había ni cuánta se estaba usando.}

\textbf{Segunda: el certificado lo firmó el presidente municipal.} Para acreditar de dónde saca
el agua una finca privada, el instrumento que se presenta ante la Provincia es una
certificación de Augusto Regis\index{Regis, Augusto}. El municipio no otorgaba derechos de agua, pero certificaba
hechos sobre el terreno, y esa certificación valía como prueba. Es la misma función que en
1938 cumplió su informe en el trámite de Campo Alegre y en 1924 su dictamen en el de Urquiza\index{Urquiza, Francisco S.}.

\textbf{Y ese rechazo no fue el final.} Tres años después, en noviembre de 1942, el mismo
solicitante obtuvo por el Decreto Nº 6705 el \textbf{reconocimiento de derecho exclusivo} de
las aguas del arroyo del Abra, con el mismo fundamento por el cual en 1939 se le había negado
la inscripción: que el arroyo nace y muere dentro de su propiedad.

Lo que en 1939 impedía registrar el derecho --- ser agua de dominio privado --- en 1942 sirve
para declararlo \textbf{exclusivo}. No es una contradicción: son dos actos distintos, uno
sobre el registro y otro sobre el derecho. Pero el resultado práctico merece anotarse.
\textbf{Un curso de agua entero quedó reconocido como de uso exclusivo de un particular, por
decreto del gobernador}, en un departamento donde ochenta años después la zona alta de uno
de sus municipios sigue esperando la red.

\textbf{Y tercera: «Abra de Lesser».} La finca fue deslindada judicialmente en 1938, según un
edicto del mismo corpus, y su nombre remite al partido de Lesser de 1909 y a la casa quinta
de diez hectáreas rematada allí en 1927. En el mismo paraje, en 2005, se iniciaría el
expediente del club de campo «Finca Lesser», que en 2014 gestionaba ciento quince metros cúbicos diarios del río Lesser con carácter permanente.

\section{Y el reparto viene de 1866}

Hay un antecedente anterior a todo lo demás de este capítulo, y aparece en el número 131 del
Boletín, en 1910, dentro de un fallo del Superior Tribunal.

\begin{ficha}{Superior Tribunal de Justicia --- B.O.\ Nº 131, 12 de febrero de 1910}
«\textit{Juicio contencioso-administrativo --- Los señores doña Gumercinda F.\ de Figueroa,
doctor Manuel Landívar y don Antonino Díaz contra el Exmo.\ Gobierno de la Provincia, sobre
\textbf{inconstitucionalidad de la Ley que dota de agua al pueblo de General Güemes}.}»

Los demandantes fundan la acción «\textit{en que la ley de 25 de Setiembre del año 1908 lesiona
sus derechos de propiedad en el caudal de agua del \textbf{Río Mojotoro}, porque al hacerse
efectiva dicha ley ese caudal de agua disminuirá las porciones que tienen adquiridas sobre ella
en virtud de una \textbf{ordenanza municipal de Campo Santo del 18 de Noviembre de 1866},
sancionada como ley por la H.\ Legislatura y promulgada por el Ejecutivo en Enero 27 del
año 67}».

El Tribunal ---doctores Ovejero, Saravia y López, sorteados--- no resuelve el fondo: resuelve
que la acción no es contencioso-administrativa sino demanda ordinaria por inconstitucionalidad,
porque lo atacado es una ley de la Legislatura y no un acto del poder administrador, y que
«\textit{las partes pueden clasificar su acción, por cualquier causa, erróneamente, sin que esto
importe alterar ó cambiar la naturaleza de la causa}».
\end{ficha}

Lo que importa para este capítulo no es el fallo sino lo que la demanda da por sabido.

\textbf{En 1866 una ordenanza municipal repartía el agua del río Mojotoro}, y la Legislatura la
sancionó como ley al año siguiente. \textbf{Diecinueve años antes de la concesión de 1885}, el
mecanismo ya estaba en uso y ya había sido convalidado por la Provincia: el municipio reparte,
la Legislatura ratifica, el derecho queda adquirido. Es exactamente la doctrina que el Decreto
14091 de 1931 pondrá por escrito y que el decreto de 1942 aplicará a la finca «Sausal o
Curuzú».

\textbf{Y en 1908 una ley provincial reasignó ese caudal a otro pueblo}, y los titulares de los
derechos de 1866 fueron a la justicia a defenderlos. \textbf{El conflicto por el agua del
Mojotoro entre quienes la tenían adquirida y un pueblo al que una ley se la da está planteado
en 1908, es decir en el primer año del Boletín.} Veintiuno después, la Ley 1269 dará a Campo
Santo potestad para vigilar y prohibir tomas desde la confluencia; veintitrés después, el
Decreto 14091 mandará investigar todas las concesiones de los dos ríos; y treinta y seis
después, la Provincia concederá cinco litros por segundo a los Ferrocarriles declarando que no
sabe si el río los lleva.

\begin{cautela}
El fallo no nombra a La Caldera. Lo que está en disputa es el caudal del Mojotoro, que se forma
con el Wierna y el de La Caldera, y ninguno de los tres demandantes aparece en el resto de este
relevamiento. \textbf{Las tres fechas normativas ---1866, 1867, 1908--- se transcriben como las
da el fallo} y no fueron contrastadas contra las normas mismas, que este libro no obtuvo.
\pendiente{ordenanza municipal de Campo Santo del 18 de noviembre de 1866, la ley que la
sanciona de enero de 1867, y la ley del 25 de septiembre de 1908 que dota de agua a General
Güemes}
\end{cautela}

\section{El municipio concedía agua en 1885}

El capítulo \ref{cap:planteo} deja abierta una pregunta: cuándo perdió el municipio la voz que tenía sobre el
agua de su jurisdicción. Un decreto de 1942 no la contesta del todo, pero corre el punto de
partida cincuenta y siete años hacia atrás y muestra el mecanismo por el cual esa facultad
cambió de manos.

\begin{ficha}{Decreto Nº 5933 del 24 de abril de 1942 --- Expediente 2462-S}
\textbf{Guillermo Schaefer}, propietario de la finca «Sausal o Curuzú», distrito de Vaqueros,
solicita al Poder Ejecutivo el \textbf{reconocimiento} de los derechos al agua del río
Wierna «otorgado por la \textbf{Municipalidad de la Caldera en el año 1885} al señor Eustaquio
Murúa, antecesor en la posesión de la finca».

Consta, por testimonio del Director del Registro Inmobiliario, que «por \textbf{acuerdo de la
Municipalidad de la Caldera de fecha 4 de junio de 1885} se concedió a don Eustaquio Murúa
el uso de \textbf{treinta litros de agua por segundo} del río Wierna» para riego de esa finca.

El decreto invoca la Ley Orgánica de Municipalidades y el Código Rural vigentes en aquella
época, y reconoce el derecho «sin perjuicio del derecho que pudieran tener otras personas al
uso de dicha agua». Ordena \textbf{inscribirlo en los libros de la Dirección General de Obras
Públicas y del Registro Inmobiliario}.
\end{ficha}

Tres cosas.

\textbf{La primera: el municipio otorgaba concesiones de agua en 1885.} No dictaminaba, no
informaba, no certificaba: \textbf{concedía}. Treinta litros por segundo de un río, por
acuerdo de su propio cuerpo deliberativo, veintitrés años antes de que existiera el Boletín
Oficial de la Provincia.

Eso amplía hacia atrás lo que el capítulo \ref{cap:expropiacion} documenta para 1914, cuando el municipio
reglamentó por ordenanza el reparto íntegro del caudal del río Santa Rufina reconociendo
«derechos adquiridos por uso tradicional». Aquellos derechos tradicionales que la ordenanza de
1914 reconocía podían ser, precisamente, concesiones municipales como ésta.

\textbf{La segunda: el Estado provincial reconoció esa facultad como legítima.} Cincuenta y
siete años después, el Poder Ejecutivo no discute que la Municipalidad de La Caldera pudiera
conceder agua en 1885: verifica el acuerdo, cita las normas que se lo permitían y ratifica el
derecho.

\textbf{Y la tercera es el mecanismo.} El acto no confirma la concesión municipal dejándola
donde estaba: la \textbf{convierte}. La inscribe en los libros de un organismo provincial y
del Registro Inmobiliario. Desde ese momento, el derecho ya no consta en un acuerdo del
municipio sino en un decreto del gobernador y en un registro de la Provincia.

Este libro no puede afirmar que ese haya sido el modo en que el municipio perdió su facultad
--- haría falta la serie completa de reconocimientos y la norma que los ordena, y ninguna de
las dos está en este relevamiento. Pero \textbf{sí puede describir lo que el decreto hace}: un
derecho otorgado por un municipio pasa, por vía de reconocimiento, a existir únicamente en
los registros provinciales. Y en 1942 no fue un caso aislado: ese mismo año el archivo
registra otros dos reconocimientos de derechos de agua en el departamento.

\section{Y en 1931 la Provincia escribió la regla}

Once años antes de ese reconocimiento hay un decreto que enuncia, con las normas en la mano,
de quién era la competencia sobre estas aguas. Es la pieza que faltaba.

\begin{ficha}{Decreto Nº 14091 --- Salta, 23 de octubre de 1931 --- Expte.\ Nº 4920-M}
Dictado por el Interventor Nacional a raíz de una nota de los representantes del Ingenio
Azucarero «La Mendieta» y del Comisionado Municipal de Campo Santo, sobre las dificultades
entre esa Municipalidad y vecinos de Mojotoro por derechos de riego.

Sus considerandos establecen: que los ríos y sus cauces son bienes públicos del Estado (arts.\
2340 inc.\ III y 2341 del Código Civil); que el régimen de concesión está en el art.\ 112 del
Código Rural; y que «\textit{perteneciendo desde un punto de vista más local, a las
atribuciones de las \textbf{Municipalidades de Campaña}, arreglar convenientemente el servicio
y distribución de las aguas de regadío de uso común en sus jurisdicciones respectivas …
estando facultadas para dictar al efecto los reglamentos y ordenanzas correspondientes,
\textbf{mientras no se sancione una ley general de Irrigación}}» (art.\ 56 inc.\ V de la Ley
Orgánica de Municipalidades).

\textbf{Art.\ 1º.} Nombra una Comisión «ad-hoc» integrada por el Director General de Obras
Públicas y Topografía, ing.\ Nolasco F.\ Cornejo, y por los Comisionados Municipales de Campo
Santo y de La Caldera, \textbf{don Pedro Mesple y don Augusto Regis}\index{Regis, Augusto},
encargada «\textit{del estudio e investigación de las concesiones otorgadas a particulares por
el Poder Ejecutivo de la Provincia para el uso-goce de las aguas de dominio público de los
\textbf{Ríos Vaqueros y Wierna}}», de las ordenanzas y reglamentos municipales que reglen su
distribución «\textit{y en general, sobre los usos y costumbres que existan sobre su
aprovechamiento}», con obligación de elevar los antecedentes y un estudio crítico.
\end{ficha}

Cuatro cosas, y ninguna es menor.

\textbf{La regla está escrita, y es la que este libro venía deduciendo.} Hasta que hubiera ley
general de irrigación, la distribución del agua de riego de uso común en cada jurisdicción era
atribución municipal, con facultad de dictar ordenanzas. Eso explica la ordenanza de 1914, la
concesión de 1885 y la competencia que el decreto de 1942 ratifica.

\textbf{Y está fechada por una condición, no por un plazo}: «mientras no se sancione una ley
general de Irrigación». El capítulo \ref{cap:aguabaja} muestra que en \textbf{septiembre de
1944} la Provincia seguía concediendo agua en precario «\textit{hasta tanto entre en vigencia
el decreto ley de agua de la Provincia}». Trece años después, la condición no se había
cumplido.

\textbf{Alguien ordenó investigar todas las concesiones del Wierna y del Vaqueros, en 1931.}
En ese momento el Wierna llevaba, por lo menos, los treinta litros por segundo de la concesión
municipal de 1885 y los cien litros por segundo que el Poder Ejecutivo había otorgado a
Francisco S.\ Urquiza\index{Urquiza, Francisco S.} en 1924 ---y \textbf{ese mismo año de 1931
la Corte de Justicia anuló la revocación de esa concesión}, como consigna el capítulo
\ref{cap:politica}. La comisión debía estudiar exactamente eso.

\textbf{Y el comisionado municipal de La Caldera que integró la comisión es Augusto Regis}, el
mismo que presidirá el municipio durante buena parte de las dos décadas siguientes y será su
comisario de policía en 1942.

\begin{cautela}
\textbf{Este libro no encontró el informe de esa comisión}, ni ningún acto posterior que lo
apruebe, lo rechace o lo cite. No sabe si llegó a reunirse. Lo que el decreto prueba es que en
octubre de 1931 el Gobierno consideró que las concesiones sobre estos dos ríos requerían un
estudio, y que puso a los dos municipios a hacerlo junto con su repartición técnica.
\pendiente{informe y antecedentes de la Comisión ad-hoc del Decreto 14091/1931 sobre las
concesiones de agua de los ríos Vaqueros y Wierna}
\end{cautela}

\textbf{Y encaja con lo que la Legislatura había hecho dos años antes.} En 1929 la Ley 1269 ---original 11078--- había dado a la Municipalidad de Campo Santo potestad para vigilar y prohibir tomas desde la confluencia del Wierna con el de La Caldera; en 1931 el Interventor pone a esa misma Municipalidad, junto con la de La Caldera, a investigar las concesiones de los dos ríos. \textbf{El arreglo institucional no es una autoridad de cuenca: son los municipios, de a dos, mirando el mismo río desde orillas distintas.}

\section{Tres puentes, y los tres viejos}

El capítulo \ref{cap:defensas} documenta que en 1935 los vecinos de La Caldera pidieron a Vialidad un puente
que uniera el pueblo con el camino nacional, y que la respuesta fue que se dispondría lo que
correspondiera «en su oportunidad». Noventa años después, el Plan de Manejo describe así la
conectividad del departamento:

\begin{ficha}{Plan de Manejo, página 113}
«Existen tres puentes; a) el que comunica la localidad de \textbf{Vaqueros con el Departamento
Capital}; b) el que atraviesa el \textbf{río Wierna} comunicando los municipios de Vaqueros y
La Caldera; c) el que permite el \textbf{acceso a La Caldera}. Se trata de infraestructuras
\textbf{antiguas, obsoletas, sin mantenimiento y muy estrechos}.»
\end{ficha}

\textbf{Tres puentes para todo el departamento, y los tres calificados de obsoletos por el
documento que el propio Estado encargó.}

El diagnóstico del Plan agrega que el principal problema del ejido urbano de La Caldera «está
relacionado a la accesibilidad ya que cuenta con \textbf{una única vía de acceso estrecha}», y
que por falta de conectividad este-oeste la mancha urbana crece hacia el sur «mediante un
camino precarizado por la falta de infraestructura hidráulica».

De modo que el corredor por el que este libro sigue la conversión de suelo tiene, en su
extremo norte, \textbf{una sola entrada}. Y el pedido de un segundo puente está hecho desde
1935.

\subsection*{Y en 2022 se licitó el puente}

La «oportunidad» llegó en marzo de 2022.

\begin{ficha}{Licitación del puente sobre el río Vaqueros y la Circunvalación Este, marzo de 2022}
Obra en la \textbf{Ruta Nº 9 --- Circunvalación Este, tramo Salta--La Caldera} ---la cobertura la llama «provincial»; la 9 es nacional---:
desde la rotonda de avenida Bolivia --- «El Quirquincho» --- hasta el \textbf{puente sobre el
río Wierna}, en el kilómetro 1.612 de la ruta nacional 9.

Cuatro kilómetros de carretera ---así en la cobertura; entre la rotonda y el puente del Wierna la distancia es bastante mayor, de modo que uno de los dos datos debe cotejarse con el pliego--- con \textbf{cuatro carriles}, tres rotondas, y un
\textbf{nuevo puente sobre el río Vaqueros} con veredas y ciclovías.

Presupuesto oficial: \textbf{más de mil trescientos cuarenta y un millones de pesos}.

\textbf{Financiada por Vialidad Nacional}; Vialidad de Salta a cargo del proceso licitatorio
y del control de obra.

Tres oferentes: Dal Borgo Construcciones S.R.L.; la unión transitoria de \textbf{DESA --- Defensa
y Encauzamiento S.A.} con Noroeste Construcciones; e Ingeniero Medina S.A.
\end{ficha}

\textbf{Primera: ochenta y siete años después del pedido de 1935 se licita un puente en el corredor, aunque no el que se pidió.} Lo que se licita: el puente de 2022 forma parte de una circunvalación de
cuatro carriles cuyo objeto declarado es dar «más seguridad vial y fluidez al tránsito» entre
la capital y el norte. \textbf{El corredor obtiene un cruce nuevo ---sobre el río Vaqueros, no el que el pueblo pidió--- como efecto de una obra pensada para pasar de largo.}

\textbf{Segunda: la paga Vialidad Nacional.} Es el mismo reparto que el capítulo \ref{cap:defensas} documenta
para 1980, cuando la Nación licitó las defensas de hormigón de la Ruta 9 mientras la Provincia
contrataba por vía administrativa las de los ríos. \textbf{Cuando lo que está en juego es la
traza nacional, aparece el presupuesto nacional.}

\textbf{Y tercera: entre los oferentes está una empresa llamada Defensa y Encauzamiento.} No
es una coincidencia de nombre: DESA es una sociedad anónima salteña, con asamblea convocada
en el Boletín Oficial de 2025, que se presentó en unión transitoria a una licitación de obra
vial sobre el río de este departamento. \textbf{El rubro que da nombre a la empresa es el que
este libro sigue desde 1910.}

\begin{cautela}
\textbf{Y hay una cuarta cosa, que este libro tenía repartida en dos capítulos y no había puesto junta.}

Los tres oferentes de la licitación del puente aparecen, con esa misma razón social o con el apellido de su titular, en el \textbf{cuadro oficial de concesiones de áridos} que la Secretaría de Minería y Energía elevó al juzgado en 2022 (capítulo \ref{cap:resistencias}).

\begin{center}\small
\begin{tabular}{@{}p{3.7cm}p{5.3cm}p{3.2cm}@{}}
\toprule
\textbf{Oferente (2022)} & \textbf{En el cuadro de Minería} & \textbf{Informe de impacto} \\
\midrule
\textbf{Noroeste Construcciones} & «Noroeste Construcciones S.A.», cantera homónima (expte.\ 19.065) & \textbf{No aprobado} \\
Dal Borgo Construcciones S.R.L. & César Domingo \textbf{Dal Borgo}: «La Vaquera» (18.770) y servidumbre de acopio (21.615) & \textbf{No aprobado} (las dos) \\
Ingeniero Medina S.A. & Héctor Enrique \textbf{Medina}: «La Generosa» (16.893) y «Tomasito» (17.049) & \textbf{No aprobado} (las dos) \\
\bottomrule
\end{tabular}
\end{center}

\textbf{Para Noroeste Construcciones S.A.\ no hay cuestión de homonimia}: es la misma razón social, literal, en las dos fuentes. Para las otras dos corresponde la reserva de siempre ---Dal Borgo Construcciones S.R.L.\ no es lo mismo que César Domingo Dal Borgo, ni Ingeniero Medina S.A.\ que Héctor Enrique Medina---, y se resuelve con dos legajos societarios que el apéndice \ref{cap:pedidos} pide.

Lo que queda registrado, y no requiere ninguna inferencia, es esto: \textbf{las tres ofertas presentadas para construir un puente sobre el río coinciden con concesionarios de extracción de áridos del mismo sistema hídrico, y los cinco derechos mineros involucrados ---cuatro canteras y una servidumbre de acopio--- figuran con la evaluación ambiental no aprobada}. Es el cruce de dos nóminas públicas ---un cuadro oficial incorporado a un expediente judicial y la cobertura de una licitación--- y cualquiera puede repetirlo.

Este libro \textbf{no afirma} irregularidad alguna, y hay una lectura enteramente inocua: en un mercado provincial chico, las firmas que tienen movimiento de suelo y áridos son las mismas que se presentan a la obra vial que los consume. Pero es exactamente lo que un vecino planteó en el taller de 2023 y que el capítulo \ref{cap:amparo} transcribe: «\textit{si la misma persona está de los dos lados del mostrador, el medio ambiente no va a mejorar}». La diferencia entre esa frase y este cuadro es que el cuadro tiene números de expediente.
\end{cautela}

\begin{cautela}
Este libro no pudo establecer si esa licitación se adjudicó ni a quién. Que la obra está en ejecución lo indica la prensa de enero de 2026, que informa que la crecida del río Vaqueros arrasó parte del nuevo puente (capítulo \ref{cap:vaqueros}). La preadjudicación estaba prevista para el 31 de marzo de 2022.

Tampoco pudo establecer si DESA participó de alguna de las obras de defensa que el capítulo \ref{cap:defensas}
documenta. \textbf{Que una empresa de ese rubro se presente a una licitación vial en esta
cuenca es un dato; no es una relación probada con ninguna otra obra del libro.}
\end{cautela}

\section{Cuántos hogares, exactamente}

Este capítulo inventaría obras y abastecimientos privados. El censo permite ponerle proporción.

\begin{ficha}{Censo Nacional 2022 --- hogares del departamento sin servicio de red}
\textbf{Sin agua por red pública:} 32,5\,\% en la fracción 66077-01, 17,6\,\% en la 02 y \textbf{58,0\,\%} en la 03.

\textbf{Sin desagüe cloacal conectado a red:} \textbf{más del 95\,\% en las tres}.

Promedio de las cincuenta fracciones de la Capital: 5,8\,\% sin agua y 13,6\,\% sin cloaca.
\end{ficha}

\begin{cautela}
\textbf{El orden de esas tres cifras es el contrario del que el resto del capítulo haría esperar.}

El municipio con gas de red, con proyecto ejecutivo de acueducto, planta potabilizadora y redes maestras aprobado en 2014, y con casi dos tercios de la población departamental ---Vaqueros, la fracción 01--- tiene \textbf{32,5\,\%} de hogares sin agua de red. El pueblo que no tiene ninguna de las tres cosas ---la fracción 02--- tiene \textbf{17,6\,\%}: \textbf{poco más de la mitad}.

No es una paradoja: es la forma exacta del crecimiento que este libro describe. El casco de La Caldera es viejo, chico y está servido; Vaqueros creció por loteos periféricos y de ladera que la red no siguió, y su promedio municipal carga con esa periferia. \textbf{La cobertura no se deteriora donde el Estado no llegó nunca: se deteriora donde el suelo se convirtió más rápido que la cañería.}

Y eso vuelve a poner en su lugar la obra de 2014. Un acueducto con redes maestras para la localidad no alcanza a los loteos que se abren después, ni a los que se abren arriba. La pregunta del apéndice \ref{cap:pedidos} ---hasta dónde llegó esa obra--- deja de ser una curiosidad de expediente: es la que explica un tercio de los hogares de Vaqueros.
\end{cautela}

El detalle y su lectura están en el capítulo \ref{cap:aguabaja}. Acá basta el contraste que da título a este capítulo: \textbf{en el departamento que provee agua a la capital, entre uno de cada seis y casi seis de cada diez hogares no tienen agua de red}, y prácticamente ninguno tiene cloaca.

\section{Septiembre de 2026: una red que sí se extiende}

Este capítulo se llama «Las redes que no están». Conviene cerrarlo con la excepción, que es de
hace días.

\begin{ficha}{Prensa, 5 de septiembre de 2026}
\textbf{«La Caldera: ampliarán la red de agua para cientos de vecinos de barrio Santiago
Apóstol».} Los trabajos comienzan en los próximos días. \textbf{Aguas del Norte} se reunirá con
los vecinos para informar plazos y tareas.
\end{ficha}

\textbf{El barrio Santiago Apóstol es el del capítulo \ref{cap:amparo}}, y la obra la ejecuta la
misma prestadora cuyo convenio sobre la Zona de Exclusión el capítulo \ref{cap:resistencias}
encuentra caducado en 2022.

\begin{cautela}
\textbf{Es una nota de prensa y este libro la consigna como tal}: acredita el anuncio, no la
obra. No tiene expediente, ni monto, ni plazo, ni cantidad de conexiones ---«cientos de
vecinos» no es una cifra---.

Pero registra algo que este capítulo debe registrar por honestidad: \textbf{la red se está
extendiendo}. Lo que el capítulo documenta es que durante décadas no se extendió, no que sea
imposible que se extienda.

\textbf{Y deja la pregunta que importa}: si el barrio Santiago Apóstol recibe red en 2026, qué
pasó entre su aprobación y hoy, y qué barrios y parajes siguen sin ella.
\end{cautela}

\section{La única autoridad metropolitana tiene fecha}

Este capítulo y el \ref{cap:aguabaja} sostienen que la provincia sabe crear figuras metropolitanas cuando
decide hacerlo, y que la única que construyó para este corredor es la del transporte de
pasajeros. Ahora esa afirmación tiene fecha y expediente.

En enero de 2002, la Secretaría de Transporte y Servicios Públicos llama a \textbf{Licitación
Pública Nacional Nº 01/02} para la «concesión del \textbf{corredor metropolitano de
transporte} regular de pasajeros entre localidades de los departamentos \textbf{Capital, La
Caldera y Cerrillos}».

El pliego cuesta \textbf{seis mil pesos} y se vende en dos lugares: el Centro Cívico Grand
Bourg y la Casa de Salta en Buenos Aires. La apertura se fija para el 28 de febrero.

Tres cosas conviene retener.

\textbf{La concesión del corredor existe desde 2002 y agrupa exactamente a los tres departamentos del aglomerado.} La autoridad vino dos años después. La Ley 7322 ---sancionada el 21 de octubre de 2004 y promulgada con veto parcial por el Decreto 2593 del 8 de noviembre de ese año--- estableció la Región Metropolitana de Salta para el transporte de pasajeros, integrada por los municipios de Salta, San Lorenzo, Vaqueros, Cerrillos, Rosario de Lerma, Campo Quijano, La Merced y La Caldera; creó la Autoridad Metropolitana del Transporte como ente autárquico, con potestades de planificación, regulación y control, y creó la sociedad prestadora SAETA. El número del decreto de promulgación lo confirman las resoluciones de la propia Autoridad publicadas en el Boletín Oficial; la fecha de sanción y la edición del Boletín (Nº 17.004) se tomaron de una transcripción del texto de la ley. De modo que en 2002 la escala metropolitana era todavía una concesión; en 2004 pasó a ser una autoridad con competencia sobre los dos municipios de este departamento. No es una construcción reciente ni improvisada: es una concesión licitada a
escala nacional, con pliego, plazos y apertura pública.

\textbf{El pliego cuesta seis mil pesos.} Póngase esa cifra al lado del encauzamiento del arroyo Urquiza que el capítulo \ref{cap:defensas} documenta para noviembre de 2001, todavía en pesos convertibles: el derecho a \textit{ofertar} por el corredor de transporte valía casi lo mismo que \textit{encauzar} un arroyo del departamento.



\begin{ficha}{Decreto Nº 4354/09 --- B.O.\ Nº 18.207, 14 de octubre de 2009}
Aprueba el Convenio de Cooperación para la creación del \textbf{Corredor Intermunicipal para el Desarrollo Sustentable (CorInDeS)}, firmado en julio de 2009 por los intendentes de La Caldera (Héctor Miguel Calabró), Vaqueros (Miguel Ángel Alemán), San Lorenzo (Ernesto Fernando Gonza), San José de los Cerrillos (Humberto Rubén Corimayo) y Salta (Miguel Ángel Isa), y por el gobernador Urtubey «como manifestación de interés». Marco: la cooperación descentralizada con la región francesa Champagne-Ardenne, declarada de interés provincial por el Decreto 2813/09 (B.O.\ Nº 18.145, 14/07/2009). Ese decreto describe el trabajo previo: un acta acuerdo de los cinco intendentes, reuniones periódicas con sus representantes técnicos y pasantes franceses que elaboraban un diagnóstico territorial turístico y propuestas de estructura intermunicipal tomando como modelo los Parques Naturales Regionales de Francia, con la reserva provincial Finca Las Costas como territorio piloto.

Los municipios se constituyen como \textbf{microrregión}. Entre los objetivos: estrategias conjuntas de gestión intermunicipal; recuperación y creación de áreas protegidas; y \textbf{«establecer los lineamientos básicos de un ordenamiento territorial, propiciando que el crecimiento urbano y las actividades productivas \ldots\ se realicen con criterios de sustentabilidad»}. Órgano: una Asamblea de Intendentes que debía reunirse al menos cada tres meses, con un representante técnico por municipio. Plazo: antes del 31 de diciembre de 2009, un protocolo adicional con el estatuto y la estructura del Corredor. Vigencia: desde la aprobación por cada municipio. El texto publicado deja en blanco el día y la ciudad de la firma.
\end{ficha}

\begin{cautela}
Esto obliga a matizar lo que este apartado sostiene. La escala metropolitana no se construyó sólo para el transporte: en 2009 los cinco municipios del corredor ---los dos de este departamento entre ellos--- se constituyeron por convenio como microrregión, con el ordenamiento territorial y las áreas protegidas entre sus fines, y la Provincia lo aprobó por decreto. Es el único instrumento que este libro encontró que trata al corredor norte como una unidad de planificación, y es anterior a todos los planes municipales que analizan los capítulos \ref{cap:pdua}, \ref{cap:cot} y \ref{cap:vaqueros}.

Lo que el libro no encontró es qué pasó después. Ni el protocolo con el estatuto, ni las ordenanzas municipales que debían ponerlo en vigencia, ni una sola resolución de la Asamblea de Intendentes, ni una mención del Corredor en el plan urbano de 2015 o en el Código de 2016. Hasta que aparezca alguna de esas piezas, el CorInDeS es una figura creada y no documentada en funcionamiento: la diferencia con la autoridad del transporte no es que una exista y la otra no, sino que una tiene ley, organismo y presupuesto, y la otra, un convenio. \pendiente{protocolo adicional con el estatuto del CorInDeS, ordenanzas de aprobación de La Caldera y Vaqueros, y resoluciones de su Asamblea de Intendentes}

Hay, sí, un rastro posterior de coordinación que no pasa por el CorInDeS. La Ordenanza 535/18 de Vaqueros declara que sus criterios para aprobar urbanizaciones fueron acordados por las autoridades de los municipios del Área Metropolitana del Valle de Lerma «en el ámbito del Comité Metropolitano de Ordenamiento Territorial», que el proyecto se había trabajado en 2016 «en conjunto entre todos los Municipios del Área Metropolitana», y que las faltas y sanciones se publicarán en un «Registro Metropolitano de Urbanizaciones» (capítulo \ref{cap:vaqueros}). La coordinación del suelo a escala metropolitana existió, entonces, al menos como acuerdo de criterios; lo que este libro no encontró es el acto que creó ese comité, ni el registro, ni si La Caldera adoptó una ordenanza equivalente. \pendiente{acto de creación del Comité Metropolitano de Ordenamiento Territorial del Área Metropolitana del Valle de Lerma, sus actas, el Registro Metropolitano de Urbanizaciones y las ordenanzas de los demás municipios que adoptaron el mismo manual}
\end{cautela}

\textbf{Y veinticuatro años después de aquella licitación no hay equivalente para la cuenca.} El río que nace en un
municipio, atraviesa otro y desemboca en un tercer departamento no tiene autoridad común; el
colectivo que hace ese mismo recorrido, sí. \textbf{La escala metropolitana existe en este
corredor para trasladar personas hacia la capital, y no para administrar el agua que baja de
él.}

\section{1946: la Provincia raciona el Wierna}

Entre el derecho de 1942 y el registro masivo de 1951 hay un acto que muestra quién decidía, en
esos años, cuánta agua bajaba del departamento. Salió en el Boletín del 22 de octubre de 1946
y el barrido del año completo lo encontró en la imagen, porque la capa de texto de la edición
lo traía cortado.

\begin{ficha}{Decreto Nº 1966 H del 16 de octubre de 1946 (en Acuerdo de Ministros)}
Expediente 8045/1946, Ministerio de Hacienda, Obras Públicas y Fomento. B.O.\ Nº 2.686, 22 de
octubre de 1946, páginas 5 y 6. Firman Lucio A.\ Cornejo, José T.\ Solá Torino y Juan W.\ Dates.

\textbf{Quién lo pide:} los regantes del río Siancas o Mojotoro, del departamento de Campo
Santo, y la \textbf{Municipalidad de Campo Santo}, que denuncia «\textit{la desigual distribución
del agua en los ríos Wierna y Mojotoro}».

\textbf{Qué constata la Dirección General de Hidráulica:} divide el sistema en tres zonas. En la
\textbf{Zona I} ---«\textit{Río Wierna, La Caldera y sus afluentes}»--- hay seis tomas que riegan
760 hectáreas y levantan 1.123 litros por segundo. En la Zona II, entre la confluencia de los
ríos La Caldera y Vaqueros y la primera toma de Campo Santo, cuatro tomas riegan 46 hectáreas
con 236 litros. En la Zona III, el sistema de Campo Santo, 7.057 hectáreas reciben 1.196 litros.
La escasez de abajo, dice el decreto, tiene por causa «\textit{el uso discrecional del agua por
los regantes de la cuenca superior}».

\textbf{Qué dispone:} mientras el Wierna lleve menos de 2.500 litros por segundo, los regantes de
las fincas \textbf{Las Mesadas, Vaqueros, Villar Urquiza, Entre Ríos, Yacones} y la de Silvano
Murúa no pueden derivar más de \textbf{382 litros por segundo} (art.\ 1); los de la Zona II, no
más de 24 (art.\ 2); los de Campo Santo, hasta 1.893 (art.\ 3). El reparto lo hace la Dirección
General de Hidráulica según un cuadro de distribución (art.\ 4). Y el artículo 5: «\textit{Las
Municipalidades de Campo Santo y de La Caldera, prestarán a la Dirección General de Hidráulica
la colaboración necesaria para la ejecución del presente decreto}», con facultad de pedir la
fuerza pública (art.\ 6).
\end{ficha}

Dos meses después, el \textbf{Decreto Nº 2599 H} del 18 de diciembre (mismo expediente; B.O.\ Nº 2.736)
reconoce los servicios de José Godoy y Delfín Chocobar como \textbf{jueces de distribución de las
aguas de los ríos Wierna y Mojotoro} durante octubre y noviembre, a ciento cincuenta pesos por
mes, pagados por la Provincia a pedido de la Dirección General de Hidráulica. Es la ejecución
del racionamiento: el reparto lo hacen agentes provinciales en el río.

\textbf{Primera: la Zona I pierde dos tercios del agua que usaba.} De 1.123 a 382 litros por
segundo, un 66 por ciento menos, mientras el río no supere los 2.500. El decreto lo funda en
el poder de policía de la Provincia sobre las aguas públicas y cita a Marienhoff: la
administración «\textit{procede por sí, sin necesidad de recurrir a la vía judicial}».

\textbf{Segunda: la voz la tuvo el municipio de abajo.} Es la continuación de lo que el apartado
sobre la Ley 1269 mostró para 1929: Campo Santo reclama, la Provincia decide. La Municipalidad
de La Caldera no aparece en los considerandos: aparece en el artículo 5, \textbf{como
colaboradora de la ejecución de un decreto que recorta el riego de su propio territorio}. Es el
mismo papel que el capítulo le encontró en 1924, 1938 y 1939 ---dictaminar, informar,
certificar---, ahora con un matiz: ya no se le pide un dato sobre el terreno, se le pide que
ayude a cumplir.

\textbf{Tercera: lo que el municipio tenía para el agua.} El balance de la comuna del último
trimestre de 1945, publicado en enero de 1946 (B.O.\ Nº 2.473, aviso Nº 1.452), registra ingresos
por despacho de bebidas, boliches, degolladura, cementerio y bailes públicos, y un solo gasto
vinculado con el agua: «\textit{Acequias compuestas}», tres pesos en octubre. El presupuesto mensual
de la comuna no llegaba a ciento treinta pesos. \textbf{Un municipio con esos recursos no
estaba en condiciones de repartir un río}, y en 1946 ya no se esperaba que lo hiciera.

\begin{cautela}
\textbf{El decreto no dice en qué departamento está cada finca del artículo 1.} Vaqueros y Yacones
son del departamento; las demás no se ubicaron. Las tres dotaciones que imprime el decreto
(1,46, 5,12 y 0,16 litros por segundo por hectárea) no son exactamente el cociente de sus
propios datos, y el artículo 6 nombra en singular «la Municipalidad mencionada» después de
nombrar dos; se transcriben como están. El cuadro de distribución de fojas 28 y el Decreto
Nº 1654 H del 23 de septiembre de 1946, que el decreto invoca como precedente para el río
Horcones, no están en el corpus. Las transcripciones provienen del barrido de 1946 (280
ediciones, Nº 2.465 a 2.744), leídas sobre la imagen porque la capa de texto del Boletín venía
cortada. \pendiente{cuadro de distribución de fs.\ 28 y antecedentes del expediente 8045/1946;
Decreto 1654 H/1946}
\end{cautela}

\section{1951: el año en que se registró el agua}

El apartado de este capítulo sobre el derecho de 1885 terminó en 1942, cuando un derecho otorgado por el
municipio en 1885 se inscribió en los registros provinciales. Nueve años después ---y cinco
después del racionamiento de 1946---, ese movimiento se vuelve masivo.

En un solo año, el Boletín publica al menos \textbf{siete actos} sobre derechos de agua del
departamento: reconocimientos de concesión sobre los ríos La Caldera y Wierna, una concesión
del Santa Rufina para una finca, la anulación y rehabilitación de la concesión de la finca
San José, y edictos citatorios para que quien se considere afectado comparezca.

Los solicitantes tienen nombre: Marcos Abraham Gutiérrez, Emilio Batel, César Pereyra Rozas,
Ana María Cornejo de Durand.

\textbf{Casi todos son \textit{reconocimientos}, no concesiones nuevas.} Es decir: derechos que
ya se ejercían y que en 1951 pasan a constar en un registro provincial, con expediente,
edicto y decreto del gobernador. Es exactamente lo que en 1942 se hizo con el derecho de
1885, ahora en escala.

Este libro \textbf{no puede afirmar} que 1951 sea el año en que la Provincia completó el
traspaso: harían falta la norma que ordenó el empadronamiento y la serie completa de
reconocimientos, y ninguna de las dos está en este relevamiento. Lo que registra es que en
ese año el departamento aparece siete veces en el registro de aguas públicas, y que \textbf{el
municipio no interviene en ninguno de los siete}.

\section{La ley que le dio el agua a otro departamento}

Antes de las redes hubo una ley, y conviene leerla porque fija un reparto que este capítulo
va a encontrar sin resolver un siglo después.

\begin{ficha}{Ley provincial Nº 1269 (original Nº 11078), sancionada el 27 de septiembre de 1929}
«Autorízase a la \textbf{Municipalidad del Departamento de Campo Santo} a efectuar los
trabajos que conceptúe necesarios en la \textbf{Playa del Río Mojotoro}, para
\textbf{reconcentrar el agua, desde las juntas del Río de Wierna con el de La Caldera},
vigilar y \textbf{prohibir se abran nuevas bocas-tomas no autorizadas por Ley}.»

Promulgada por Cornejo el 1º de octubre de 1929, con la firma del ministro L.\ C.\ Uriburu, y publicada en el Boletín Oficial Nº 1.292 del 11 de octubre. No fija monto ni plazo.\par\textit{En el Digesto provincial figura como \textbf{Ley 1269 (original Nº 11078)}: el número original es el que imprime el propio Boletín de 1929 y el que consigna la fuente que el Digesto cita, la \textit{Recopilación General de Leyes} compilada por Gavino Ojeda. No es una anomalía: en esos años leyes y decretos se registraban en una misma numeración correlativa ---la ley de defensas de Campo Santo de 1928 es la 9707, y el decreto que la ejecuta, el 9718---. \textbf{Ya no rige}: el Digesto de la Cámara de Diputados la lista entre las leyes de objeto cumplido, plazo vencido o no vigentes. Firmaron la sanción W.\ Toledo, C.\ Arias Aranda, E.\ Campilongo y Marcos B.\ Figueroa.}
\end{ficha}

Léase el alcance. La Legislatura autoriza a un municipio —Campo Santo— a intervenir sobre
el agua \textbf{desde el punto donde el río Wierna se junta con el de La Caldera}, es
decir, sobre el caudal que baja del departamento vecino. Y le da además dos potestades que
no son de obra sino de policía: \textbf{vigilar}, y \textbf{prohibir nuevas tomas}.

Es de 1929 y es la contracara exacta del episodio que el capítulo \ref{cap:politica} reconstruye: un año
antes, los vecinos y la comisión municipal de Campo Santo habían reclamado contra una
concesión de agua otorgada aguas arriba, en La Caldera, y habían perdido en abril para
ganar en mayo. Al año siguiente obtuvieron por ley lo que el reclamo administrativo no les
había dado de manera estable: autoridad sobre el punto donde el agua entra a su
jurisdicción.

Lo que interesa a este capítulo es lo que esa ley revela sobre la forma de resolver el
problema. La cuenca es una y atraviesa varios departamentos; el arreglo institucional
consiste en darle a uno de ellos potestades sobre el tramo que le importa. No se crea una
autoridad de cuenca: se reparte el río.

Un siglo después, el capítulo \ref{cap:vaqueros} documenta que el río Vaqueros separa dos municipios y que
ninguno administra el conjunto del cauce; que la extracción de áridos se concesiona en
ambas márgenes sin autoridad común; y que la única figura metropolitana que la provincia
construyó para el corredor —con planificación, regulación y poder de policía— es la del
transporte de pasajeros.

\section{Qué significa que una concesión sea eventual}

El apartado anterior mostró que hay agua sobre la que el Estado no lleva registro. Falta la
otra distinción, la que separa a las concesiones entre sí, y que este libro leyó mal hasta
obtener el texto del Código de Aguas.

\begin{ficha}{Una concesión de agua --- Boletín Oficial Nº 19.330, 26 de junio de 2014}
Gestión ante la Secretaría de Recursos Hídricos, publicada por cinco días conforme los arts.\ 24 inc.\ a), 51, 62, 65, 66, 140/158 y 201 del Código de Aguas y su Decreto Reglamentario 2299/03.

\textbf{Gestionante:} Abel Oscar Almeida (h), titular registral del dominio del \textbf{Catastro Nº 1700 del departamento La Caldera}.

\textbf{Objeto:} concesión de uso de \textbf{agua pública subterránea de pozo perforado}, para \textbf{consumo humano}, para \textbf{abastecer a 120 habitantes}.

\textbf{Volumen:} 0,00876 hm\textsuperscript{3} anuales --- 1,60 m\textsuperscript{3} por hora con quince horas de bombeo diarias.

\textbf{Carácter: eventual} (art.\ 143 del Código de Aguas).

Expediente Nº 0090034-25.217/2014-0/3. Publicado cinco días consecutivos, del 25 de junio al 1 de julio de 2014, en los Boletines Nº 19.329 a 19.333. Plazo de oposición: treinta días hábiles desde la última publicación, para quienes tengan derecho o interés legítimo, conforme el art.\ 309 del Código de Aguas.

Firma: Dra.\ Silvia F.\ Santamaría\index{Santamaría, Silvia F.}, Abogada Jefe del Programa Fiscalización y Control. Orden de publicación 100042106.
\end{ficha}

\begin{cautela}
La aritmética del aviso cierra y conviene hacerla explícita, porque describe una política de hecho. Un metro cúbico y sesenta por hora durante quince horas son veinticuatro metros cúbicos diarios; multiplicados por trescientos sesenta y cinco días dan los 8.760 metros cúbicos anuales que el aviso expresa en hectómetros. Repartidos entre ciento veinte habitantes, son \textbf{exactamente doscientos litros por persona y por día}. La dotación está calculada con precisión, para consumo humano permanente.

Y sin embargo la concesión se pide, y se publica, con \textbf{carácter eventual}. Obtenido el texto del Código de Aguas, esa aparente anomalía se disuelve y deja algo mejor en su lugar.
\end{cautela}

\begin{ficha}{Art.\ 143 de la Ley 7017 --- texto}
``Carácter de las concesiones. Las concesiones de uso de aguas subterráneas, salvo las que se hubieran obtenido por declaración judicial de aguas privadas, \textbf{serán eventuales}.''
\end{ficha}

\begin{cautela}
Corresponde una rectificación. Este apartado había leído el carácter eventual de la concesión de 2014 como una asimetría entre la forma jurídica y la función real --- un abastecimiento domiciliario permanente encuadrado como uso precario. \textbf{No lo es}: el artículo 143 establece que \textbf{todas} las concesiones de agua subterránea son eventuales, sin más excepción que las declaradas judicialmente privadas. No hubo elección de encuadre; hubo aplicación de la regla general.

Pero la regla general dice, ella misma, algo más fuerte que la anomalía que este libro creyó ver. Los titulares de concesiones eventuales, en el esquema del Código, reciben su dotación \textbf{después} de satisfechas las permanentes y según el orden de otorgamiento. De modo que \textbf{en la provincia de Salta nadie que dependa de un pozo tiene derecho firme al agua que extrae}: por definición legal, su título cede ante cualquier concesión permanente y ante la disponibilidad del recurso.

Eso reordena lo que este capítulo viene describiendo. La infraestructura que se resuelve lote por lote --- perforación propia, concesión eventual --- no es sólo más cara y más desigual que la red: es \textbf{jurídicamente más frágil}. Ciento veinte personas abastecidas por un pozo en 2014 lo están bajo un título que la ley califica de eventual, mientras quien recibe agua de red la recibe bajo concesión permanente. La diferencia entre tener red y no tenerla no se mide sólo en cañerías.
\end{cautela}

\begin{cautela}
Lo estructural es más simple. En 2014, en un departamento sin red de agua suficiente, la manera de abastecer a ciento veinte personas era que un particular perforara un pozo en su propia parcela y pidiera una concesión. Eso es lo que este capítulo describe cuando dice que la infraestructura se resuelve lote por lote: no es una metáfora, es un trámite con número de orden de publicación.
\end{cautela}

Y no es la única del año. En agosto de 2014, \textbf{María Eugenia Carante}, propietaria del catastro 2907 del departamento, gestionó la asignación de riego que registraba el catastro de origen Nº 126 a favor de la única matrícula rural resultante ---y en 2021 volvería a gestionar que ese mismo riego pasara del catastro 2907 a uno nuevo, el 6582: el derecho siguió a la tierra de subdivisión en subdivisión---, y solicitó concesión de agua pública \textbf{para irrigación de diez hectáreas, con carácter permanente}, con una dotación de 5,25 litros por segundo, a derivar del \textbf{río Wierna, margen derecha}. Expediente 0090034-168.839/14, firma de la asesora legal Sandra Mabel Siegrist.

Dos observaciones.

\textbf{La primera:} esta concesión es \textbf{permanente}, no eventual. La diferencia con la del Catastro 1700 no está en el destino sino en la fuente: aquélla era de agua subterránea de pozo, que el artículo 143 declara eventual sin excepción; ésta es de agua superficial derivada de un río, que admite el carácter permanente. La precariedad, entonces, no alcanza a todos por igual: alcanza a quien depende del subsuelo. Aunque, en los hechos, no a todos los que dependen de él: en 2015, dos gestiones para loteos sobre agua subterránea ---el dren de El Durazno y el pozo de La Misión III--- se publicaron con carácter permanente (capítulo \ref{cap:loteo}). La regla del artículo 143 alcanzó a los ciento veinte vecinos del Catastro 1700 y no a esos dos emprendimientos.

\textbf{La segunda:} el trámite consiste en trasladar una asignación de riego preexistente ``a favor de la única matrícula rural resultante''. Es decir, hubo una subdivisión, y el derecho de agua que correspondía al inmueble de origen se concentra en la fracción que quedó como rural. \textbf{El agua sigue al suelo cuando el suelo se fracciona}, y lo hace por un trámite de tres días de publicación. Es la contracara exacta de lo que este capítulo describe: mientras la red pública no llega, los derechos de agua preexistentes se reasignan al ritmo de la subdivisión de la tierra.

\section{El círculo}

Los capítulos anteriores permiten cerrar el razonamiento.

El municipio recibe una norma de calidad, financiada y asistida por la Provincia y el CFI, que presupone un aparato técnico. Su departamento absorbe en doce años un crecimiento poblacional del 58,4 por ciento, el cuarto del país, sin red de gas en su territorio ---en el departamento sólo la tiene Vaqueros---, con cobertura cloacal parcial y con desigualdad documentada en el acceso al agua. No tiene ese aparato ni los recursos para sostenerlo, porque su base imponible crece en superficie loteada y no en recaudación. La superficie loteada crece porque la norma, cuando el suelo no es apto, no prohíbe sino que deriva a un procedimiento agravado que el municipio tampoco tiene capacidad de evaluar. Y cuando el resultado se inunda, la reparación recae sobre el mismo municipio, que la financia con la misma base.

Este circuito no requiere mala fe de nadie para funcionar exactamente así, y ése es su rasgo definitorio. No hay una decisión de capturar el Estado local: hay un Estado local al que se le asignó una función sin los medios para ejercerla, y un mercado que ocupa el espacio que esa incapacidad deja libre.

Lo que sí es una decisión ---provincial, sostenida y verificable--- es la de transferir la competencia sin transferir la capacidad. Ahí hay responsabilidad política, y no es municipal.

